% =====================================================================
% 本編『Git & GitHub 入門 — はじめてのバージョン管理』
% エンジン: LuaLaTeX（jlreq / book）
% ビルド:   latexmk -lualatex main_intro.tex
%
% 対象: Git / GitHub をはじめて触る人。1人で使い、GitHubに上げるまで。
% チーム開発（ブランチ・PR・コンフリクト）は別冊 main_team.tex を参照。
% =====================================================================
\documentclass[book, paper=a4, fontsize=10.5pt, titlepage]{jlreq}

% =====================================================================
% preamble.tex — パッケージ読込・色・カスタム環境・マクロの定義
% 本文（chapters/*.tex）から使う見た目の部品はすべてここに集約する。
% エンジン: LuaLaTeX + jlreq
% =====================================================================

% --- 基本パッケージ ---------------------------------------------------
\usepackage{graphicx}          % 画像の取り込み
\usepackage{xcolor}            % 色指定
\usepackage{booktabs}          % 罫線のきれいな表
\renewcommand{\arraystretch}{1.25} % 表の行を少しゆったりさせる
\usepackage{tabularx}          % 幅を自動調整する表
\usepackage{enumitem}          % 箇条書きの調整
\usepackage{fancyvrb}          % コード枠に使う Verbatim
\usepackage[most]{tcolorbox}   % 装飾枠（コード・注意書き）
\usepackage{adjustbox}         % 図が本文幅を超えたときの自動縮小
\usepackage{tikz}              % 概念図・フロー図
\usetikzlibrary{arrows.meta,positioning,shapes.geometric,calc,fit,backgrounds}
\usepackage{float}             % 図を本文の位置に固定する [H]
\usepackage{caption}           % 図番号つきキャプションの体裁
\captionsetup{font=small, labelfont=bf, labelsep=quad}

% 画像の探索パス（既存の Markdown 版と同じ images/ を共有する）
\graphicspath{{../images/}{./images/}}

% --- 配色 -------------------------------------------------------------
% Markdown 版の Mermaid 図の色をそのまま踏襲する（薄い色は使わない方針）。
\definecolor{gitRed}{HTML}{C0392B}     % 作業ツリー・注意
\definecolor{gitOrange}{HTML}{D68910}  % ステージング
\definecolor{gitGreen}{HTML}{2E7D32}   % リポジトリ・完了
\definecolor{gitBlue}{HTML}{4A90D9}    % コミット・情報
\definecolor{gitPurple}{HTML}{6C3483}  % 自分側
\definecolor{gitTeal}{HTML}{117A65}    % 両方採用
\definecolor{gitAmber}{HTML}{8A6D0B}   % 判断ノード
\definecolor{codeBg}{HTML}{F4F5F7}     % コード枠の背景
\definecolor{codeFrame}{HTML}{D0D3D9}  % コード枠の枠線

% --- コード枠 ---------------------------------------------------------
% 使い方: \begin{codebox}[bash] ... Verbatim ... \end{codebox}
% 省略時のラベルは bash。右上に言語ラベルを小さく表示する。
\newtcolorbox{codebox}[1][bash]{%
  enhanced, breakable,
  colback=codeBg, colframe=codeFrame, boxrule=0.6pt, arc=2pt,
  left=9pt, right=9pt, top=13pt, bottom=8pt,
  before skip=8pt, after skip=8pt,
  overlay={%
    \node[anchor=north east, font=\sffamily\scriptsize\bfseries,
          text=gray!60] at ([xshift=-6pt, yshift=-3pt]frame.north east) {#1};%
  },
}
% コード本体を書くための Verbatim 環境（小さめ・行間やや詰め）
\DefineVerbatimEnvironment{code}{Verbatim}{fontsize=\small, xleftmargin=0pt}

% 本文中のインラインコード: \code{git status}
\newcommand{\cmd}[1]{\texttt{#1}}

% --- 注意書き（admonition）枠 ----------------------------------------
% ヒント（緑）・注意（赤）・補足（グレー）の3種。emoji は使わずラベルで表す。
% タイトルは本文（明るい背景）に直接乗るため、白文字だと見えない。
% タイトル文字色は各テーマ色にして、はっきり読めるようにする。
\newtcolorbox{hintbox}{%
  enhanced, breakable,
  colback=gitGreen!5, colframe=gitGreen, boxrule=0.8pt, arc=2pt,
  left=8pt, right=8pt, top=6pt, bottom=6pt,
  coltitle=gitGreen, fonttitle=\sffamily\bfseries\small,
  title=ヒント, attach title to upper=\quad,
}
\newtcolorbox{warnbox}{%
  enhanced, breakable,
  colback=gitRed!5, colframe=gitRed, boxrule=0.8pt, arc=2pt,
  left=8pt, right=8pt, top=6pt, bottom=6pt,
  coltitle=gitRed, fonttitle=\sffamily\bfseries\small,
  title=注意, attach title to upper=\quad,
}
% 執筆メモ（雛形用の枠）。原稿を清書したら、この枠ごと削除する。
\newtcolorbox{draftbox}{%
  enhanced, breakable,
  colback=gitAmber!6, colframe=white, boxrule=0pt, arc=2pt,
  borderline={0.9pt}{0pt}{gitAmber, dashed},
  left=8pt, right=8pt, top=6pt, bottom=6pt,
  coltitle=gitAmber, fonttitle=\sffamily\bfseries\small,
  title=執筆メモ, attach title to upper=\quad,
}
\newtcolorbox{notebox}{%
  enhanced, breakable,
  colback=gray!8, colframe=gray!60!black, boxrule=0.8pt, arc=2pt,
  left=8pt, right=8pt, top=6pt, bottom=6pt,
  coltitle=gray!30!black, fonttitle=\sffamily\bfseries\small,
  title=補足, attach title to upper=\quad,
}

% --- スクリーンショット取り込み --------------------------------------
% \screenshot[幅の割合]{ファイル}{キャプション}
% フロートにせず本文の流れに沿って中央寄せで置く（説明の直後に見せたいため）。
\newcommand{\screenshot}[3][0.7]{%
  \begin{center}
    \includegraphics[width=#1\linewidth]{#2}\\[3pt]
    {\small\itshape #3}%
  \end{center}%
}

% \screenshottrim[幅の割合]{ファイル}{キャプション}{trim指定}
% 余白の多いスクリーンショットを切り詰めて大きく見せる。
% trim は「左 下 右 上」の順で、画像1ピクセル = 1bp として指定する
% （例: 1920x1080 の下半分を捨てる → {0 540bp 0 0}）。
\newcommand{\screenshottrim}[4][0.95]{%
  \begin{center}
    \includegraphics[width=#1\linewidth, trim=#4, clip]{#2}\\[3pt]
    {\small\itshape #3}%
  \end{center}%
}

% --- TikZ 共通スタイル ------------------------------------------------
% Mermaid のフロー図を再現するためのノード／矢印スタイル。
\tikzset{
  flownode/.style={%
    rounded corners=4pt, draw=none, text=white,
    font=\sffamily\small, align=center,
    minimum height=1.15cm, minimum width=2.2cm, inner sep=6pt},
  decisionnode/.style={%
    diamond, aspect=1.9, draw=none, text=white,
    font=\sffamily\small, align=center, inner sep=3pt},
  flowarrow/.style={-{Latex[length=3.6mm, width=3mm]}, line width=1.8pt, draw=black},
  edgelabel/.style={font=\sffamily\normalsize\bfseries, fill=white, inner sep=3pt},
}
% 図全体を中央寄せで置くためのヘルパ環境（図番号なし）。
% 本文幅を超えた図は自動で縮小する（幅が足りている図は拡大しない）。
\newenvironment{diagram}%
  {\par\medskip\begin{center}\begin{adjustbox}{max width=\linewidth}}%
  {\end{adjustbox}\end{center}\par\medskip}

% --- 図番号つきの図（本文が「図1.2のように」と参照する概念図・フロー図用）---
% 手順に紐づく操作スクショには使わない（そちらは \screenshot を使う）。
\makeatletter
% TikZ など本体を囲む環境: \begin{numfig}{キャプション}{ラベル} ... \end{numfig}
\newenvironment{numfig}[2]{%
  \def\nf@cap{#1}\def\nf@lab{#2}%
  \begin{figure}[H]\centering\begin{adjustbox}{max width=\linewidth}%
}{%
  \end{adjustbox}\caption{\nf@cap}\label{\nf@lab}\end{figure}%
}
\makeatother
% 画像1枚を図番号つきで置く: \figimage[幅]{ファイル}{キャプション}{ラベル}
\newcommand{\figimage}[4][0.95]{%
  \begin{figure}[H]\centering
    \includegraphics[width=#1\linewidth]{#2}
    \caption{#3}\label{#4}
  \end{figure}%
}


% hyperref はプリアンブルの最後（他パッケージ読込後）に置くのが安全。
\usepackage[unicode, hidelinks]{hyperref}

% --- 書誌情報 ---------------------------------------------------------
\title{Git \& GitHub 入門}
\author{}
\date{\today}

\begin{document}

% --- 表紙 -------------------------------------------------------------
\begin{titlepage}
  \centering
  \vspace*{4cm}
  {\sffamily\bfseries\Huge Git \& GitHub 入門\par}
  \vspace{1.2cm}
  {\sffamily\Large はじめてのバージョン管理\par}
  \vfill
  {\large \today\par}
\end{titlepage}

% --- この本のねらい（表紙の次のページ） ------------------------------
\thispagestyle{empty}
\vspace*{1.5cm}
{\sffamily\LARGE\bfseries はじめに —この本を読むとできること—\par}
\vspace{1em}

本書は、\textbf{Git・GitHub をはじめて使う人}のための入門書です。特別な前提知識は要りません。読み終えると、\textbf{自分のPCと GitHub アカウントだけ}で、次のことができるようになります。

\begin{itemize}[leftmargin=1.6em]
  \item Git で、ファイルの変更を「コミット」として記録し、履歴を残せる
  \item 変更の履歴や差分を確認できる（\cmd{git status} / \cmd{git log} / \cmd{git diff}）
  \item \cmd{.gitignore} で、記録したくないファイルを除外できる
  \item GitHub アカウントを作り、\cmd{gh} でログイン（認証）できる
  \item 手元のリポジトリを GitHub に上げ（\cmd{push}）、最新を取り込める（\cmd{pull} / \cmd{fetch}）
  \item \textbf{「\cmd{pull} → 編集 → \cmd{add} → \cmd{commit} → \cmd{push}」} の日常の流れを回せる
\end{itemize}

つまり、\textbf{1人で Git を使い、その成果を GitHub にバックアップ・共有する}ところまでが、本書のゴールです。

\begin{hintbox}
\textbf{この先へ}（本書を読み終えたら）
\begin{itemize}[leftmargin=1.4em, nosep]
  \item チームでブランチや PR を使って開発したい → 別冊『チームで安全に開発する』
  \item コマンドをすぐ引きたい → 『逆引きチートシート』
  \item エラーや操作ミスで困った → 『トラブルシューティング』
\end{itemize}
\end{hintbox}

\clearpage

% --- 目次 -------------------------------------------------------------
\tableofcontents
\clearpage

% --- 本文 -------------------------------------------------------------
% 本編は「教科書」に徹する。トラブルシューティング・チートシートは
% 別冊 main_team.tex の付録にまとめている。
% =====================================================================
% 第1章 Git編 — 自分のPCの中でGitを使う
% 元原稿: 01_Git編.md
% =====================================================================
\chapter{Git編 —自分のPCの中でGitを使う—}

この章では、\textbf{自分のPCの中だけで完結するGitの基本}を説明します。GitHub（クラウド）とつなぐ話は第2章で扱います。まずはここで「変更を記録する」感覚をつかみましょう。

\begin{hintbox}
\textbf{この章の読み方}
\begin{itemize}[leftmargin=1.4em, nosep]
  \item とりあえず動かしたい人 → 「インストール」→「最初の設定」→「【実践】はじめてのコミット」だけ写経すればOK
  \item コマンドをド忘れした人 → 『逆引きチートシート』へ
  \item エラーが出て困った人 → 『トラブルシューティング』へ
\end{itemize}
\end{hintbox}

% ---------------------------------------------------------------------
\section{そもそもGitとは}

Gitは\textbf{ファイルの変更履歴を記録・管理するツール}です。一言でいうと「\textbf{いつでも戻れるセーブポイントを作れる仕組み}」です。

たとえば、レポートをこんなふうに管理していませんか。

\begin{codebox}[text]
\begin{code}
レポート.docx
レポート_最終.docx
レポート_最終2.docx
レポート_本当に最終.docx
\end{code}
\end{codebox}

このやり方だと、次のことに困ります。

\begin{itemize}[leftmargin=1.4em]
  \item \textbf{どれが本当に最新か分からない}
  \item \textbf{ファイルが増え続けて散らかる}
  \item \textbf{いつ・誰が・どこを・なぜ変えたのか分からない}
\end{itemize}

コピーを取っておけば「昔の状態に戻す」ことはできます。でも「\textbf{どこが違うのか}」「\textbf{なぜ変えたのか}」までは分かりません。ここがGitとの決定的な違いです。

Gitを使うと、\textbf{ファイルは1つのまま}、変更の節目ごとに「\textbf{コミット}」という記録（セーブポイント）を残せます。各コミットには「いつ・誰が・何を・なぜ変えたか」に加えて、\textbf{前回との差分}まで残るので、後からいつでも正確にたどれます。

\begin{numfig}{コミットを積み重ねると、変更履歴が1本の線としてつながる}{fig:commit-chain}
\begin{tikzpicture}[node distance=1.1cm]
  \node[flownode, fill=gitBlue] (A) {昔のレポート\\（コミットA）};
  \node[flownode, fill=gitBlue, right=of A] (B) {少し進めた\\（コミットB）};
  \node[flownode, fill=gitGreen, right=of B] (C) {完成\\（コミットC）};
  \draw[flowarrow] (A) -- (B);
  \draw[flowarrow] (B) -- (C);
\end{tikzpicture}
\end{numfig}

% ---------------------------------------------------------------------
\section{Gitの3つのエリア（最重要の概念）}

Gitでは、ファイルは次の3つのエリアを通って記録されます。\textbf{ここだけは必ず押さえてください。}

\begin{hintbox}
下の図に出てくる \cmd{git add} と \cmd{git commit} は、この後の「実践」で実際に使うコマンドです。ここでは\textbf{エリアの間を移す操作の名前}とだけ捉えてください。
\end{hintbox}

3つのエリアは、\textbf{机 → 段ボール箱 → 倉庫}にたとえると分かりやすいです（図\ref{fig:areas-concept}）。作業ツリーは原本を広げた\textbf{作業机}、ステージングは変更を詰める\textbf{段ボール箱}、ローカルリポジトリは箱を保管する\textbf{倉庫}です。

\figimage[0.95]{01_git/git_area.png}{3つのエリアの概念図（机 → 段ボール箱 → 倉庫）}{fig:areas-concept}

この流れをコマンドで表すと、\cmd{git add} と \cmd{git commit} で次のエリアへ送っていくことになります（図\ref{fig:areas-flow}）。

\begin{numfig}{3つのエリアと、それらをつなぐ \texttt{git add} / \texttt{git commit}}{fig:areas-flow}
\begin{tikzpicture}
  \node[flownode, fill=gitRed, text width=2.4cm] (W)
        {作業ツリー\\{\scriptsize (Working Tree)}\\原本を広げた作業机};
  \node[flownode, fill=gitOrange, text width=2.4cm, right=3.6cm of W] (S)
        {ステージング\\{\scriptsize (Staging Area)}\\写しを詰める段ボール箱};
  \node[flownode, fill=gitGreen, text width=2.4cm, right=3.6cm of S] (R)
        {ローカルリポジトリ\\{\scriptsize (Local Repository)}\\写しを保管する倉庫};
  \draw[flowarrow] (W) -- (S) node[edgelabel, midway, above] {git add};
  \draw[flowarrow] (S) -- (R) node[edgelabel, midway, above] {git commit};
\end{tikzpicture}
\end{numfig}

\begin{center}
\begin{tabularx}{\linewidth}{@{}l X X@{}}
\toprule
エリア & これは何か & たとえると \\
\midrule
\textbf{作業ツリー} & 実際にファイルを編集している、いつものフォルダ & 原本を広げた作業机 \\
\textbf{ステージング} & 「次のコミットに含める変更」を詰めるエリア & コミット1回分の段ボール箱 \\
\textbf{ローカルリポジトリ} & コミットとして履歴が保存されるエリア（\cmd{.git} の中） & 箱を保管する倉庫 \\
\bottomrule
\end{tabularx}
\end{center}

大事なのは、\cmd{git add} も \cmd{git commit} も「\textbf{写しを取って次のエリアに置く}」操作だという点です。元のエリアには元のまま残るので、先に進めても空になりません。

\begin{itemize}[leftmargin=1.4em]
  \item \cmd{git add}：机の上の原本の写しを箱に入れる（原本は机に残る）
  \item \cmd{git commit}：箱の中身の写しを倉庫に保管する（箱の中身も残る）
\end{itemize}

\begin{notebox}
Gitは各エリアが「その時点の写し（スナップショット）」を持つように動きます。だからコミットは\textbf{差分ではなくスナップショット}として記録されます。
\end{notebox}

つまり基本操作は \textbf{「編集する → \cmd{add} → \cmd{commit}」} の繰り返しです。本書ではこれを\textbf{黄金フロー}と呼びます。この3ステップが分かれば、Gitの8割は理解できたようなものです。

% ---------------------------------------------------------------------
\section{Gitのインストール}

まず、Gitが\textbf{すでに入っていないか}を確認しましょう。ターミナル（WindowsはPowerShellやコマンドプロンプト）で次を実行します。

\begin{codebox}[bash]
\begin{code}
git --version
# 出力例: git version 2.45.0
\end{code}
\end{codebox}

\begin{hintbox}
上のように\textbf{バージョンが表示されたら、Gitはすでにインストール済み}です。以下のOS別のインストール手順は\textbf{飛ばして}、次の「最初の設定」へ進んでください。\cmd{command not found}（見つからない）と出た場合だけ、お使いのOSの手順に従ってインストールします。インストール後にもう一度 \cmd{git --version} を実行し、バージョンが表示されれば成功です。
\end{hintbox}

\begin{warnbox}
\textbf{インストール直後は、開いていたターミナルではまだ \cmd{git --version} が反映されないことがあります。} PATH（コマンドの置き場所を示す設定）がターミナルの起動時に読み込まれるためです。\cmd{command not found} のままなら、\textbf{ターミナルをいったん閉じて開き直して}から、もう一度実行してください。
\end{warnbox}

\subsection{Windows}

Windowsでは「Git for Windows」をインストールします。これには \textbf{Git Bash} という、macOS/Linux風のコマンドが使えるターミナルが付いてきます。本書のWindowsの説明はこの Git Bash を使う前提です。

\begin{enumerate}[leftmargin=1.6em]
  \item \url{https://git-scm.com/download/win} からインストーラをダウンロードし、ダブルクリックで起動します。
  \item インストーラの画面が続きます。\textbf{基本はすべてデフォルトのまま「Next」で進めて問題ありません。} ただし次の3つの画面だけは選択を確認してください。
\end{enumerate}

\noindent\textbf{① 既定のエディタ}

デフォルトはVimですが不慣れだと難しいため、VS Codeを使う予定なら \textbf{「Use Visual Studio Code as Git's default editor」} を選びます。
\screenshot[0.62]{01_git/win_install_05.png}{既定エディタの選択}

\noindent\textbf{② 初期ブランチ名}

\textbf{「Override the default branch name…」} を選び、\cmd{main} と入力します（世界的に \cmd{main} が標準）。
\screenshot[0.62]{01_git/win_install_06.png}{初期ブランチ名の設定}

\noindent\textbf{③ 改行コードの扱い}

WindowsとほかのOSで改行コードが異なるため重要です。複数環境で開発するなら \textbf{「Checkout as-is, commit Unix-style line endings」} を選びます。
\screenshot[0.62]{01_git/win_install_10.png}{改行コードの設定}

残りの画面はデフォルトのまま進め、最後に「Install」をクリックします。インストール後、スタートメニューから \textbf{Git Bash} を起動し、\cmd{git --version} でバージョンが出れば完了です。

\subsection{macOS}

次のどちらか一方でOKです。

\noindent\textbf{方法A: Xcode Command Line Tools（最も手軽）}

ターミナル（\cmd{アプリケーション > ユーティリティ > ターミナル.app}）を開き、次を実行します。

\begin{codebox}[bash]
\begin{code}
xcode-select --install
\end{code}
\end{codebox}

ダイアログが出たら「インストール」をクリックします。

\noindent\textbf{方法B: Homebrew（最新版を使いたい場合）}

Homebrew（\url{https://brew.sh/index_ja}）が入っていれば、次で最新のGitが入ります。

\begin{codebox}[bash]
\begin{code}
brew install git
\end{code}
\end{codebox}

どちらの場合も、\cmd{git --version} でバージョンが表示されれば完了です。

\subsection{Linux}

お使いのディストリビューションのパッケージマネージャを使います。自分の環境に合うものを1つ選んで実行します。

\begin{codebox}[bash]
\begin{code}
# Debian / Ubuntu 系
sudo apt update && sudo apt install git

# Fedora / RHEL 系
sudo dnf install git

# Arch 系
sudo pacman -S git
\end{code}
\end{codebox}

\cmd{git --version} でバージョンが表示されれば完了です。

% ---------------------------------------------------------------------
\section{最初の設定}

インストール直後に\textbf{一度だけ}行う設定です。コミットに記録される「あなたの名前とメールアドレス」を登録します。ここで設定した名前・メールが、後でGitHub上の履歴にも表示されます。

ターミナル（Windowsは Git Bash）を開き、次を順に実行します。

\begin{codebox}[bash]
\begin{code}
# あなたの名前を登録（GitHubのアカウント名に合わせるのがおすすめ）
git config --global user.name "your-name"

# あなたのメールアドレスを登録（GitHubで使う予定のメール）
git config --global user.email "you@example.com"

# 設定内容を確認する
git config --global --list
\end{code}
\end{codebox}

\begin{warnbox}
名前やメールにスペースが含まれる場合は、上のように \textbf{ダブルクォート \cmd{"} で囲んでください}。
\end{warnbox}

\cmd{--global} は「このPC全体の共通設定」という意味です。一度やれば、以降すべてのリポジトリに適用されます。

\begin{hintbox}
おすすめの追加設定（新規リポジトリの初期ブランチ名を \cmd{main} に統一）
\begin{codebox}[bash]
\begin{code}
git config --global init.defaultBranch main
\end{code}
\end{codebox}
\end{hintbox}

% ---------------------------------------------------------------------
\section{【実践】はじめてのコミット}

ここが本番です。\textbf{黄金フロー（編集 → \cmd{add} → \cmd{commit}）}を、実際に手を動かして体験します。\textbf{この通りに写経すれば、最初のコミットが作れます。}

\subsection{ステップ0：作業フォルダを用意して移動する}

\begin{codebox}[bash]
\begin{code}
# 練習用フォルダを作って中に入る
mkdir git-practice
cd git-practice
\end{code}
\end{codebox}

\begin{hintbox}
\cmd{cd}（change directory）はフォルダを移動するコマンドです。以降の操作は「今いるフォルダ」に対して行われます。
\end{hintbox}

\subsection{ステップ1：リポジトリを作る（git init）}

\begin{codebox}[bash]
\begin{code}
git init
\end{code}
\end{codebox}

これで今のフォルダに \cmd{.git} という隠しフォルダが作られ、Gitでの管理が始まります。\textbf{これが、あなたのPCの中の「ローカルリポジトリ」}です。履歴（コミット）はこの \cmd{.git} の中に保存されます（中身を直接触る必要はありません）。

\begin{codebox}[bash]
\begin{code}
# git initに成功すると、次のように表示されます
hoge@n171h100 git-practice % git init
Initialized empty Git repository in /Users/hoge/git-practice/.git/
\end{code}
\end{codebox}

\screenshot[0.62]{01_git/01_git_init_folder.png}{git init後の.gitフォルダ（隠しファイル表示をONにすると見えます）}

\subsection{ステップ2：ファイルを作る}

試しにテキストファイルを1つ作ります。

\begin{codebox}[bash]
\begin{code}
echo "Hello Git" > hello.txt
\end{code}
\end{codebox}

\subsection{ステップ3：変更を確認する（git status）}

\begin{codebox}[bash]
\begin{code}
git status
\end{code}
\end{codebox}

\cmd{hello.txt} が「Untracked files（まだ管理されていないファイル）」として赤色で表示されます。
\screenshot[0.62]{01_git/01_git_status.png}{git statusの出力}

\subsection{ステップ4：ステージングに追加する（git add）}

\begin{codebox}[bash]
\begin{code}
git add hello.txt
\end{code}
\end{codebox}

これで \cmd{hello.txt} が「次のコミットに含める変更」として選ばれました。もう一度 \cmd{git status} を実行すると、今度は緑色で表示されます。
\screenshot[0.62]{01_git/01_git_add_status.png}{git add後のgit statusの出力}

\subsection{ステップ5：コミットする（git commit）}

\begin{codebox}[bash]
\begin{code}
git commit -m "最初のコミット"
\end{code}
\end{codebox}

\cmd{-m} の後ろの文字列が\textbf{コミットメッセージ}（変更内容の一言メモ）です。これでコミット（セーブポイント）が1つ作られました。

\subsection{ステップ6：履歴を確認する（git log）}

\begin{codebox}[bash]
\begin{code}
git log
\end{code}
\end{codebox}

\screenshot[0.62]{01_git/01_git_commit_log.png}{git logの出力}

今作ったコミットが、作成者・日時・メッセージとともに表示されます。\cmd{q} キーで表示を終了します。

\subsection{黄金フローのまとめ}

上の\textbf{ステップ4（\cmd{git add}）→ ステップ5（\cmd{git commit}）}が黄金フローの中核です。図にすると、図\ref{fig:golden-flow}のように3つの繰り返しになります。

\begin{numfig}{黄金フロー：編集 → \texttt{git add} → \texttt{git commit} の繰り返し}{fig:golden-flow}
\begin{tikzpicture}[scale=1.5, transform shape]
  \node[flownode, fill=gitRed] (E) {① ファイルを編集};
  \node[flownode, fill=gitOrange, right=3.4cm of E] (S) {ステージング};
  \node[flownode, fill=gitGreen, right=3.4cm of S] (R) {コミットとして\\記録};
  \draw[flowarrow] (E) -- (S) node[edgelabel, midway, above] {② git add};
  \draw[flowarrow] (S) -- (R) node[edgelabel, midway, above] {③ git commit};
  \draw[flowarrow] (R) to[out=-30, in=-150]
        node[edgelabel, midway, below] {繰り返す} (E);
\end{tikzpicture}
\end{numfig}

実際の作業は、この \textbf{「編集 → add → commit」} をひたすら繰り返すだけです。

\begin{codebox}[bash]
\begin{code}
# 日常的に使う最小セット
git add .                       # 変更した全ファイルをステージングに追加
git commit -m "変更内容のメモ"   # コミットとして記録
\end{code}
\end{codebox}

\begin{hintbox}
\cmd{git add .} の \cmd{.}（ドット）は「今いるフォルダ以下のすべての変更」という意味です。
\end{hintbox}

\subsection{いつ \texttt{add} / \texttt{commit} すればいい？}

最初に迷うのが\textbf{コミットするタイミング}です。目安はシンプルで、\textbf{「1つの作業に区切りがついたら」コミットする}——これだけです。たとえば次のような節目です。

\begin{itemize}[leftmargin=1.4em]
  \item 機能を1つ作り終えた
  \item 文章を1節書けた／図を1つ追加した
  \item バグを1つ直した
  \item 設定を変えて、動くことを確認できた
\end{itemize}

ポイントは、\textbf{後から見て「何をしたか」を一言で言えるまとまり}にすることです。その一言が、そのまま\textbf{コミットメッセージ}（例：\cmd{ログイン画面を追加}）になります。

\begin{hintbox}
\textbf{こまめに、小さくコミットするのがおすすめ}です。小さく区切っておくと、\textbf{後で戻したいときにその1コミットだけ取り消せ}、\textbf{不具合が入った箇所も見つけやすく}なります。逆に何日分もためて1回でコミットすると、後から分けられず、メッセージも「いろいろ修正」のように曖昧になりがちです。
\end{hintbox}

\textbf{\cmd{add} と \cmd{commit} の役割の違い}も、このタイミングで押さえておきましょう。

\begin{itemize}[leftmargin=1.4em]
  \item \cmd{git add}：これから作るコミットに\textbf{含める変更を選ぶ}（コミット直前の準備）
  \item \cmd{git commit}：選んだ変更で\textbf{セーブポイントを確定する}（区切りをつける）
\end{itemize}

ふだんは \cmd{git add .} → \cmd{git commit} を\textbf{続けて実行すればOK}です。「変更した中の一部だけをコミットしたい」ときにだけ、\cmd{add} で\textbf{選ぶ}意味が出てきます。

% ---------------------------------------------------------------------
\section{状態を確認するコマンド}

Gitは「今どういう状態か」をこまめに確認しながら使うのが安全です。困ったらまず \cmd{git status} を打つ癖をつけましょう。

\begin{center}
\begin{tabularx}{\linewidth}{@{}l X@{}}
\toprule
コマンド & 何が分かるか \\
\midrule
\cmd{git status} & どのファイルが変更され、どのエリアにあるか \\
\cmd{git log} & これまでのコミット履歴（作成者・日時・メッセージ） \\
\cmd{git log -{}-oneline} & コミット履歴を1行ずつ簡潔に表示 \\
\cmd{git diff} & まだステージングしていない変更の中身（どこを書き換えたか） \\
\cmd{git diff -{}-staged} & ステージング済みの変更の中身 \\
\bottomrule
\end{tabularx}
\end{center}

このうち \textbf{\cmd{git diff}} は「どこをどう書き換えたか」を確認するのによく使います。追加された行が \cmd{+} 付きの緑、削除された行が \cmd{-} 付きの赤で表示されます。

\begin{notebox}
\cmd{git add} でステージングした後は \cmd{git diff} に何も出なくなります。そのときは \cmd{git diff -{}-staged} を使います。さらに \cmd{git commit} した後は、\cmd{git log} で確認します。
\end{notebox}
\screenshot[0.92]{01_git/01_git_log_diff.png}{git logとgit diffの出力}

% ---------------------------------------------------------------------
\section{.gitignore：記録したくないファイルを除外する}

プロジェクトには「Gitで記録したくないファイル」があります。たとえば次のようなものです。

\begin{itemize}[leftmargin=1.4em]
  \item パスワードやAPIキーなどの\textbf{秘密情報}（誤って公開すると危険）
  \item 自動生成される一時ファイル（\cmd{.DS\_Store}、\cmd{\_\_pycache\_\_/}、LaTeXの \cmd{.aux} など。作り直せるので記録する意味がない）
  \item サイズの大きいデータやビルド成果物（リポジトリが重くなる。GitHubには容量の上限もある）
\end{itemize}

これらは \cmd{.gitignore} というファイルに書いておくと、Gitが\textbf{最初から無視}してくれます（\cmd{git status} にも出てきません）。リポジトリの一番上の階層に \cmd{.gitignore} を作り、無視したいものを1行ずつ書きます。

\begin{codebox}[text]
\begin{code}
# macOS が作る隠しファイル
.DS_Store

# Python のキャッシュ
__pycache__/
*.pyc

# LaTeX の中間ファイル
*.aux
*.log
*.out

# 秘密情報（絶対にコミットしない）
.env
secrets.txt
\end{code}
\end{codebox}

\begin{warnbox}
\textbf{一度コミットしてしまったファイルは、後から \cmd{.gitignore} に書いても無視されません。} 秘密情報は\textbf{最初のコミット前}に \cmd{.gitignore} へ登録してください。
\end{warnbox}

\begin{hintbox}
言語やフレームワークごとの \cmd{.gitignore} テンプレートは \url{https://github.com/github/gitignore} で、自動生成は \url{https://www.toptal.com/developers/gitignore} で手に入ります。
\end{hintbox}

% ---------------------------------------------------------------------
\section*{この章のまとめ}
\addcontentsline{toc}{section}{この章のまとめ}

\begin{itemize}[leftmargin=1.4em]
  \item Gitは「いつでも戻れるセーブポイント」を作る変更履歴管理ツール
  \item ファイルは \textbf{作業ツリー → ステージング → ローカルリポジトリ} の3エリアへ「写し」が渡っていく（元のエリアは空にならない）
  \item 基本操作は \textbf{「編集 → \cmd{git add} → \cmd{git commit}」} の繰り返し
  \item 困ったら \cmd{git status} で現在地を確認する
\end{itemize}

次は、ここで作った履歴を\textbf{クラウド（GitHub）に上げて共有する}方法です（第2章）。

% =====================================================================
% 第2章 GitHub編 — ローカルとクラウドをつなぐ
% 元原稿: 02_GitHub編.md
% =====================================================================
\chapter{GitHub編 —ローカルとクラウドをつなぐ—}

本章は第1章の続きです。第1章では「自分のPCの中だけ」で変更を記録しました。本章では、その記録を \textbf{GitHub（クラウド）} に置いて、バックアップ・複数PC・複数人での共有ができるようにします。

扱うのは \textbf{「ローカルとGitHubをつなぎ、同期するところまで」}です。ブランチを分けた共同作業・プルリクエスト（PR）・コンフリクトの解決は第3章で扱います。

\begin{hintbox}
\textbf{この章の読み方}
\begin{itemize}[leftmargin=1.4em, nosep]
  \item まず一度つないでみたい人 → 「アカウントを作る」→「認証を設定する」→「リモートリポジトリを用意する」→「つなぐ」を順に写経
  \item コマンドをド忘れした人 → 第5章「逆引きチートシート」へ
  \item エラーが出て困った人 → 第4章「トラブルシューティング編」へ
\end{itemize}
\end{hintbox}

% ---------------------------------------------------------------------
\section{そもそもGitHubとは}

\textbf{Git と GitHub は別物です。}

\begin{itemize}[leftmargin=1.4em]
  \item \textbf{Git} … あなたのPCの中で動く「変更履歴を記録する道具」（第1章で使ったもの）
  \item \textbf{GitHub} … Gitのリポジトリを\textbf{クラウド上に置いて共有できるWebサービス}
\end{itemize}

第1章では、あなたのPCの中に「倉庫（ローカルリポジトリ）」を作りました。GitHubは、その\textbf{クラウド上にあるもう一つの倉庫（リモートリポジトリ）}だと考えてください。手元の倉庫の\textbf{写しをクラウドの倉庫にも置いておく}イメージです（図\ref{fig:local-remote}）。

\begin{numfig}{ローカルとリモートは「同期」でつながる2つの倉庫}{fig:local-remote}
\begin{tikzpicture}
  \node[flownode, fill=gitRed, text width=3.2cm] (L)
        {あなたのPC\\ローカルリポジトリ\\{\scriptsize（手元の倉庫）}};
  \node[flownode, fill=gitBlue, text width=3.2cm, right=4.5cm of L] (R)
        {GitHub\\リモートリポジトリ\\{\scriptsize（クラウドの倉庫）}};
  \draw[{Latex[length=3.6mm,width=3mm]}-{Latex[length=3.6mm,width=3mm]},
        line width=1.8pt, draw=black]
        (L) -- (R) node[edgelabel, midway, above] {同期};
\end{tikzpicture}
\end{numfig}

クラウドに置くと、こんなことができます。

\begin{itemize}[leftmargin=1.4em]
  \item \textbf{バックアップ}：PCが壊れても、GitHubから取り戻せる
  \item \textbf{複数PCで共有}：研究室のPCと自宅のPCで同じ作業を続けられる
  \item \textbf{複数人で共有}：チームで同じリポジトリを編集できる
  \item \textbf{公開}：コードや手引書を世界に公開できる（このドキュメントもGitHub上にあります）
\end{itemize}

\begin{hintbox}
「リモート（remote）」は\textbf{手元から離れた場所＝クラウド側}、「ローカル（local）」は\textbf{あなたのPCの中}を指す言葉です。以降ずっと出てくるので覚えておいてください。
\end{hintbox}

% ---------------------------------------------------------------------
\section{ローカルとリモートの全体像}

ローカルとリモートの間で写しをやり取りするコマンドは、主に次の4つです。\textbf{向き（アップロードかダウンロードか）}で整理すると分かりやすいです（図\ref{fig:remote-overview}）。

\begin{numfig}{ローカルとリモートをつなぐ4つのコマンドと向き}{fig:remote-overview}
\begin{tikzpicture}
  \node[flownode, fill=gitRed, text width=3cm, minimum height=2.2cm] (L)
        {ローカルリポジトリ\\{\scriptsize（あなたのPC）}};
  \node[flownode, fill=gitBlue, text width=3cm, minimum height=2.2cm, right=7cm of L] (R)
        {リモートリポジトリ\\{\scriptsize（GitHub）}};
  \draw[flowarrow] ([yshift=7mm]L.east) -- ([yshift=7mm]R.west)
        node[edgelabel, midway, above] {git push};
  \draw[flowarrow] ([yshift=0mm]R.west) -- ([yshift=0mm]L.east)
        node[edgelabel, midway, above] {git pull / git fetch};
  \draw[flowarrow, dashed] ([yshift=-7mm]R.west) -- ([yshift=-7mm]L.east)
        node[edgelabel, midway, below] {git clone（最初の1回）};
\end{tikzpicture}
\end{numfig}

\begin{center}
\begin{tabularx}{\linewidth}{@{}l l X@{}}
\toprule
コマンド & 向き & 何をするか \\
\midrule
\cmd{git clone} & リモート → ローカル & リモートの倉庫を\textbf{丸ごと手元に複製}する（最初の1回だけ） \\
\cmd{git push} & ローカル → リモート & 手元のコミットを\textbf{クラウドへアップロード}する \\
\cmd{git fetch} & リモート → ローカル & クラウドの最新を\textbf{取得するだけ}（手元の作業には反映しない） \\
\cmd{git pull} & リモート → ローカル & クラウドの最新を\textbf{取得して手元に反映}する（\cmd{fetch} + 統合） \\
\bottomrule
\end{tabularx}
\end{center}

\begin{hintbox}
\cmd{git add} / \cmd{git commit}（第1章）は\textbf{ローカルの中だけ}の操作でした。GitHubと通信するのは、この \cmd{clone} / \cmd{push} / \cmd{fetch} / \cmd{pull} のときだけです。
\end{hintbox}

% ---------------------------------------------------------------------
\section{GitHubアカウントを作る}

まだアカウントがなければ作成します。

\begin{enumerate}[leftmargin=1.6em]
  \item \url{https://github.com/} を開き、\textbf{「Sign up」} をクリックします。
  \item メールアドレス・パスワード・ユーザー名を入力して進めます。
  \item 登録したメールアドレスに届く確認コードを入力すれば完了です。
\end{enumerate}

\begin{hintbox}
ここで登録する\textbf{メールアドレスは、第1章の \cmd{git config -{}-global user.email} で設定したものに揃える}のがおすすめです。揃えておくと、コミットとGitHubアカウントが正しく結びつきます。
\end{hintbox}

% ---------------------------------------------------------------------
\section{認証を設定する（HTTPS と SSH）}

GitHubに \cmd{push}（アップロード）するときは、「本当にあなたか」を確認する\textbf{認証}が必要です。方式は主に2つあり、\textbf{どちらか一方}を設定すればOKです。

\begin{center}
\begin{tabularx}{\linewidth}{@{}l l X@{}}
\toprule
方式 & URLの形 & 特徴 \\
\midrule
\textbf{HTTPS} & \cmd{https://github.com/\ldots.git} & \textbf{GitHub CLI（\cmd{gh}）でログインするだけ}で使えて最も手軽。トークンの手動管理も可能 \\
\textbf{SSH鍵} & \cmd{git@github.com:\ldots.git} & 鍵を1度登録すれば以降は認証不要。複数リポジトリで使い回せる \\
\bottomrule
\end{tabularx}
\end{center}

\begin{hintbox}
迷ったら、まずは手軽な \textbf{HTTPS} で始めて構いません。HTTPSは、後述の \textbf{GitHub CLI（\cmd{gh}）でログインするだけ}で使えるようになり、トークンを自分で発行・管理する必要がありません。頻繁に使うようになったら SSH に移行しても快適です。
\end{hintbox}

\subsection{方式A：HTTPS}

HTTPSでの認証には2通りあります。\textbf{GitHub CLI（\cmd{gh}）を使う方法が圧倒的に手軽}なので、こちらを最初におすすめします。トークンを手作業で扱いたい場合のために、従来のPAT方式も後に載せます。

\subsubsection{A-1（おすすめ）：GitHub CLI（gh）でログインする}

\textbf{GitHub CLI（\cmd{gh}）} はGitHubを操作する公式コマンドです。これで一度ログインしておくと、\cmd{gh} がGitの認証を裏で肩代わりしてくれるため、\textbf{トークンを自分で発行しなくても \cmd{push} / \cmd{pull} がそのまま通ります}。

\noindent\textbf{手順1：\cmd{gh} をインストールする}

\begin{warnbox}
\textbf{\cmd{gh} はGit本体には含まれていません。} 第1章でGitをインストールしても \cmd{gh} は入っていないので、\textbf{最初に別途インストールが必要}です。まず \cmd{gh -{}-version} を実行し、\cmd{command not found}（未インストール）なら下記でインストールしてください。
\end{warnbox}

\begin{codebox}[bash]
\begin{code}
# macOS（Homebrew）
brew install gh

# Windows（winget）
winget install --id GitHub.cli

# Linux（Debian / Ubuntu 系）
sudo apt install gh
\end{code}
\end{codebox}

インストール後、次でバージョンが表示されれば成功です。

\begin{codebox}[bash]
\begin{code}
gh --version
# 出力例: gh version 2.62.0 (2024-11-27)
\end{code}
\end{codebox}

\begin{hintbox}
他の環境やうまくいかない場合は、公式手順 \url{https://cli.github.com/} を参照してください。
\end{hintbox}

\noindent\textbf{手順2：ログインする}

\begin{codebox}[bash]
\begin{code}
gh auth login
\end{code}
\end{codebox}

対話形式で聞かれるので、次のように選びます。

\begin{hintbox}
質問文や順序、選択肢の文言は \textbf{\cmd{gh} のバージョンによって変わることがあります}。下記の画面例と完全に一致しなくても、\textbf{同じ意味の選択肢}（GitHub.com／HTTPS／ブラウザでログイン など）を選べば問題ありません。
\end{hintbox}

\noindent\textbf{(1) 認証するアカウントの選択}
\begin{codebox}[text]
\begin{code}
# GitHub.com を選択する
? Where do you use GitHub?  [Use arrows to move, type to filter]
> GitHub.com
  Other
\end{code}
\end{codebox}

\noindent\textbf{(2) 認証方式の選択}
\begin{codebox}[text]
\begin{code}
# HTTPS を選択する
? What is your preferred protocol for Git operations?
> HTTPS
  SSH
\end{code}
\end{codebox}

\noindent\textbf{(3) Gitの認証をGitHubの資格情報で行うか}
\begin{codebox}[text]
\begin{code}
# Y を入力する
Authenticate Git with your GitHub credentials? (Y/n) Y
\end{code}
\end{codebox}

\noindent\textbf{(4) 認証方法の選択}
\begin{codebox}[text]
\begin{code}
# Login with a web browser を選択する
? How would you like to authenticate?
> Login with a web browser
  Paste an authentication token

# 選ぶと、次のようにワンタイムコードが表示される。
# Enter を押すとブラウザが開き、GitHubにログインして認証できる。
! First copy your one-time code: FDAC-FFFB
Press Enter to open https://github.com/login/device in your browser...
\end{code}
\end{codebox}

Enter押下の後に開くブラウザ画面で \textbf{Continue} を押します。
\screenshot[0.62]{02_github/gh_auth_login.png}{GitHubログイン画面（Continueを押す）}

続けて表示される画面で、ワンタイムコードを入力して \textbf{Continue} を押します。
\screenshot[0.62]{02_github/gh_authorize_window.png}{GitHub認証画面（ワンタイムコードを入力）}

これで完了です。\cmd{gh} がGitの認証情報ヘルパーとして登録され、\textbf{以降 \cmd{git push} / \cmd{git pull} で認証を聞かれなくなります}。トークンの発行・保管は \cmd{gh} が肩代わりしてくれるため、自分で発行・管理する必要はありません。

\begin{warnbox}
ただし \cmd{gh} が\textbf{自動で無期限に更新してくれるわけではありません}。トークンの有効期限が切れたり、GitHub側でアクセスを取り消したり、新しい権限（スコープ）が必要になった場合は、再度 \cmd{gh auth login}（権限追加なら \cmd{gh auth refresh}）で認証し直してください。\cmd{gh auth status} で現在の状態を確認できます。
\end{warnbox}

\begin{codebox}[bash]
\begin{code}
# ログイン状態を確認する
gh auth status
\end{code}
\end{codebox}

\begin{hintbox}
この方法なら、次項の「トークンを発行する」手順は不要です。\textbf{基本はこの \cmd{gh auth login} だけで済みます。}
\end{hintbox}

\subsubsection{A-2（手動）：パーソナルアクセストークン（PAT）を使う}

\cmd{gh} を使わずHTTPSで \cmd{push} すると、ユーザー名とパスワードを聞かれます。ここで、\textbf{GitHubのログインパスワードは使えません}。代わりに \textbf{パーソナルアクセストークン（PAT）} という専用の文字列を発行して使います。

\noindent\textbf{トークンの発行手順}

\begin{enumerate}[leftmargin=1.6em]
  \item GitHubにログインし、右上のアイコン → \textbf{Settings} を開きます。
  \item 左メニュー最下部の \textbf{Developer settings} をクリックします。
  \item \textbf{Personal access tokens} → \textbf{Tokens (classic)} → \textbf{Generate new token (classic)} を選びます。
  \item 次を設定して発行します。
    \begin{itemize}[leftmargin=1.4em, nosep]
      \item \textbf{Note}：用途のメモ（例：\cmd{my-laptop}）
      \item \textbf{Expiration}：有効期限（例：90 days）
      \item \textbf{Select scopes}：\cmd{repo} にチェック（リポジトリの読み書きに必要）
    \end{itemize}
  \item \textbf{Generate token} を押すと、\verb|ghp_xxxx…| という文字列が表示されます。
\end{enumerate}

\screenshot[0.62]{02_github/create_token.png}{トークン発行画面}
\screenshot[0.62]{02_github/create_token_setting.png}{作成するトークンの設定}
\screenshot[0.62]{02_github/pat_scopes.png}{トークン発行画面（repoスコープにチェック）}

\begin{warnbox}
\textbf{トークンは発行直後の1回しか表示されません。} その場でコピーし、パスワードマネージャ等に保管してください。画面を閉じると二度と確認できず、再発行が必要になります。
\end{warnbox}

\noindent\textbf{トークンの使い方}

\cmd{push} 時にユーザー名とパスワードを聞かれたら、\textbf{パスワードの欄に、先ほどのトークンを貼り付けます}。

\begin{codebox}[text]
\begin{code}
Username: your-github-name
Password: ghp_xxxxxxxxxxxxxxxxxxxx   ← ログインパスワードではなくトークンを入力
\end{code}
\end{codebox}

一度入力すればトークンが記憶されるので、次回以降は入力不要になります。

\subsection{方式B：SSH鍵}

SSHは、\textbf{鍵ペア（秘密鍵と公開鍵）}を使う方式です。公開鍵をGitHubに登録しておくと、以降は認証情報の入力なしで \cmd{push} できます。

\noindent\textbf{手順1：鍵ペアを作る}

ターミナル（WindowsはGit Bash）で次を実行します。メールアドレスはGitHubのものにします。

\begin{codebox}[bash]
\begin{code}
ssh-keygen -t ed25519 -C "you@example.com"
\end{code}
\end{codebox}

\begin{itemize}[leftmargin=1.4em]
  \item 保存先を聞かれたら、そのまま \textbf{Enter}（デフォルトの \verb|~/.ssh/id_ed25519| に保存されます）
  \item パスフレーズは、設定してもしなくても構いません（設定するとより安全）
\end{itemize}

これで秘密鍵 \verb|id_ed25519| と公開鍵 \verb|id_ed25519.pub| の2つが \verb|~/.ssh/| に作られます。

\begin{warnbox}
\textbf{秘密鍵（\cmd{.pub} が付かない方）は絶対に他人に渡さないでください。} GitHubに登録するのは公開鍵（\cmd{.pub} が付く方）だけです。
\end{warnbox}

\noindent\textbf{手順2：公開鍵の中身をコピーする}

\begin{codebox}[bash]
\begin{code}
# macOS
pbcopy < ~/.ssh/id_ed25519.pub

# Windows（Git Bash）
cat ~/.ssh/id_ed25519.pub | clip

# Linux（表示された内容を手動でコピー）
cat ~/.ssh/id_ed25519.pub
\end{code}
\end{codebox}

\noindent\textbf{手順3：GitHubに公開鍵を登録する}

\begin{enumerate}[leftmargin=1.6em]
  \item GitHubの \textbf{Settings} → \textbf{SSH and GPG keys} を開きます。
  \item \textbf{New SSH key} をクリックします。
  \item \textbf{Title} に分かりやすい名前（例：\cmd{my-laptop}）、\textbf{Key} に手順2でコピーした内容を貼り付けます。
  \item \textbf{Add SSH key} で登録完了です。
\end{enumerate}

\noindent\textbf{手順4：接続を確認する}

\begin{codebox}[bash]
\begin{code}
ssh -T git@github.com
\end{code}
\end{codebox}

次のように表示されれば成功です（\cmd{yes} を求められたら \cmd{yes} と入力）。

\begin{codebox}[text]
\begin{code}
Hi your-github-name! You've successfully authenticated, but GitHub does not provide shell access.
\end{code}
\end{codebox}

\screenshot[0.62]{02_github/ssh_key_add.png}{SSH and GPG keys の New SSH key 登録画面}

% ---------------------------------------------------------------------
\section{リモートリポジトリを用意する}

GitHub上に、写しを置くための空の倉庫（リモートリポジトリ）を作ります。

\begin{enumerate}[leftmargin=1.6em]
  \item GitHub右上の \textbf{「+」} → \textbf{New repository} をクリックします。
  \item 次を設定します。
    \begin{itemize}[leftmargin=1.4em, nosep]
      \item \textbf{Repository name}：リポジトリ名（例：\cmd{git-practice}）
      \item \textbf{Public / Private}：公開するか非公開か（研究用途なら通常 \textbf{Private}）
      \item \textbf{Add a README file}：\textbf{チェックしない}（後述のパターンBで既存ローカルをつなぐ場合、付けると手順が増えるため）
    \end{itemize}
  \item \textbf{Create repository} をクリックします。
\end{enumerate}

作成後の画面に、\cmd{https://github.com/\ldots} や \cmd{git@github.com:\ldots} で始まる\textbf{リポジトリのURL}が表示されます。次の手順で使うのでこの画面を開いたままにしておきます。

\screenshot[0.62]{02_github/new_repo.png}{新しいリポジトリの作成}
\screenshot[0.62]{02_github/new_repo_setting.png}{作成するリポジトリの設定}
\screenshot[0.62]{02_github/new_repo_create.png}{リポジトリの作成}

% ---------------------------------------------------------------------
\section{ローカルとリモートをつなぐ}

つなぎ方は、状況によって2パターンあります（図\ref{fig:connect-patterns}）。

\begin{numfig}{状況に応じた2つのつなぎ方}{fig:connect-patterns}
\begin{tikzpicture}[node distance=1.2cm and 1.6cm]
  \node[decisionnode, fill=gitAmber, text width=3.4cm] (Q)
        {これから始める？\\既にローカルがある？};
  \node[flownode, fill=gitGreen, below left=1.4cm and 0.2cm of Q, text width=3.4cm] (A)
        {パターンA\\\texttt{git clone}};
  \node[flownode, fill=gitGreen, below right=1.4cm and 0.2cm of Q, text width=3.6cm] (B)
        {パターンB\\\texttt{git remote add + push}};
  \draw[flowarrow] (Q) -- (A) node[edgelabel, midway, left] {GitHubから始める};
  \draw[flowarrow] (Q) -- (B) node[edgelabel, midway, right] {手元のを上げる};
\end{tikzpicture}
\end{numfig}

\subsection{つなぐ前に：リポジトリのURL（リンクパス）を取得する}

どちらのパターンでも、\cmd{git clone} や \cmd{git remote add} に渡す\textbf{リポジトリのURL}が必要です。このURLは\textbf{GitHubのリポジトリ画面から取得}します。取得場所は、リポジトリが空かどうかで少し変わります。

\noindent\textbf{ケース1：中身のあるリポジトリ（\cmd{clone} する場合など）}

リポジトリのトップ画面にある緑色の \textbf{「＜＞ Code」} ボタンをクリックします。開いたパネルで \textbf{HTTPS / SSH} のタブを選ぶと、それぞれのURLが表示されます。右側のコピーアイコンで丸ごとコピーできます。

\begin{itemize}[leftmargin=1.4em, nosep]
  \item \textbf{HTTPS} タブ → \cmd{https://github.com/\ldots.git}
  \item \textbf{SSH} タブ → \cmd{git@github.com:\ldots.git}
\end{itemize}

\screenshot[0.62]{02_github/repo_https_url.png}{Codeボタンから HTTPS のURLを取得する}
\screenshot[0.62]{02_github/repo_ssh_url.png}{Codeボタンから SSH のURLを取得する}

\noindent\textbf{ケース2：作りたてで空のリポジトリ（パターンBでつなぐ場合）}

README を付けずに作った直後は、\textbf{「Quick setup」} の画面が表示されます。ここに表示される \textbf{HTTPS / SSH} の切り替えボタンと、その隣のURL（コピーアイコン付き）から取得します。\cmd{\ldots or push an existing repository from the command line} の欄には、パターンBでそのまま使えるコマンドも載っています。

\screenshot[0.62]{02_github/repo_quick_setup.png}{空リポジトリのQuick setup画面とURL}

\begin{hintbox}
\textbf{HTTPS と SSH のどちらのURLを使うかは、「認証を設定する」で設定した方式に合わせます。} \cmd{gh}／トークンで設定したなら HTTPS のURL、SSH鍵を登録したなら SSH のURLを選んでください。
\end{hintbox}

\subsection{パターンA：GitHubで作って \texttt{git clone}（これから始める人向け）}

GitHub上のリポジトリを、\textbf{丸ごと手元に複製}します。作業したいフォルダの\textbf{親フォルダ}で実行してください（\cmd{clone} するとリポジトリ名のフォルダが新しく作られます）。

\begin{codebox}[bash]
\begin{code}
# HTTPS の場合
git clone https://github.com/your-name/git-practice.git

# SSH の場合
git clone git@github.com:your-name/git-practice.git
\end{code}
\end{codebox}

複製が終わると \cmd{git-practice} フォルダが作られ、その中に入れば作業を始められます。\textbf{\cmd{origin}（リモートの登録）も自動で設定される}ので、以降はすぐ \cmd{push} / \cmd{pull} できます。

\begin{codebox}[bash]
\begin{code}
cd git-practice
git remote -v   # origin が登録されていることを確認
\end{code}
\end{codebox}

次のように \cmd{origin} が \cmd{fetch} / \cmd{push} の2行で表示されれば、正しく登録されています（URLは自分のものになります）。

\begin{codebox}[text]
\begin{code}
origin  https://github.com/your-name/git-practice.git (fetch)
origin  https://github.com/your-name/git-practice.git (push)
\end{code}
\end{codebox}

\subsection{パターンB：既存のローカルを \texttt{push}（第1章で作ったものを上げる）}

第1章で \cmd{git init} して作った手元のリポジトリを、後からGitHubにつなぐ場合です。リポジトリのフォルダ内で実行します。

\begin{codebox}[bash]
\begin{code}
# 1. リモート（origin）を登録する
git remote add origin https://github.com/your-name/git-practice.git
#    SSH なら： git remote add origin git@github.com:your-name/git-practice.git

# 2. 登録内容を確認する
git remote -v
#    → origin  https://github.com/your-name/git-practice.git (fetch)
#      origin  https://github.com/your-name/git-practice.git (push)
#      のように2行表示されれば登録成功

# 3. 現在のブランチ名を main に統一する
git branch -M main

# 4. 初回のアップロード（-u で「送り先」を覚えさせる）
git push -u origin main
\end{code}
\end{codebox}

\begin{itemize}[leftmargin=1.4em]
  \item \textbf{\cmd{origin}} は「リモートリポジトリの標準的な別名」です。以降 URL の代わりに \cmd{origin} と書けます。
  \item \textbf{\cmd{-u}（\cmd{-{}-set-upstream}）} は、\cmd{push} の引数で指定した送り先（ここでは \cmd{origin} の \cmd{main}）を、いまのローカルブランチの「既定の送り先」として覚えさせるオプションです。
    \begin{itemize}[leftmargin=1.4em, nosep]
      \item 本来 \cmd{git push} は毎回どこへ送るか（\cmd{origin main}）を指定する必要があります。\cmd{-u} を一度付けておくと対応付けが記録され、\textbf{次回からは \cmd{git push} / \cmd{git pull} だけ}で同じ相手とやり取りできます。
      \item 記録された送り先は \cmd{git branch -vv} で確認できます（ブランチ名の横に \cmd{[origin/main]} と表示されます）。
    \end{itemize}
\end{itemize}

\begin{warnbox}
\cmd{git remote add origin \ldots} で \cmd{error: remote origin already exists.} と出たら、既に登録済みです。URLを付け替えたいときは \cmd{git remote set-url origin <新しいURL>} を使います。
\end{warnbox}

% ---------------------------------------------------------------------
\section{アップロードする（push）}

\cmd{git push} は、\textbf{ローカルリポジトリのコミット（写し）を、リモートへ送る}操作です。第1章の「編集 → add → commit」でローカルに記録した内容を、クラウドへ反映します（図\ref{fig:push-flow}）。

\begin{numfig}{編集からクラウド反映までの流れ}{fig:push-flow}
\begin{tikzpicture}
  \node[flownode, fill=gitRed] (E) {編集};
  \node[flownode, fill=gitOrange, right=2.2cm of E] (S) {ステージング};
  \node[flownode, fill=gitGreen, right=2.4cm of S] (L) {ローカル\\リポジトリ};
  \node[flownode, fill=gitBlue, right=2.4cm of L] (R) {GitHub\\リモート};
  \draw[flowarrow] (E) -- (S) node[edgelabel, midway, above] {git add};
  \draw[flowarrow] (S) -- (L) node[edgelabel, midway, above] {git commit};
  \draw[flowarrow] (L) -- (R) node[edgelabel, midway, above] {git push};
\end{tikzpicture}
\end{numfig}

\begin{codebox}[bash]
\begin{code}
# 初回（-u で送り先を覚えさせる）
git push -u origin main

# 2回目以降はこれだけでOK
git push
\end{code}
\end{codebox}

\cmd{push} が終わったら、GitHubのリポジトリ画面を再読み込みしてみてください。コミットしたファイルがクラウド上に表示されているはずです。

\begin{warnbox}
\textbf{コミットしていない変更は \cmd{push} されません。} \cmd{push} は「ローカルリポジトリにあるコミット」を送るものです。手元の編集をクラウドに反映したいときは、必ず先に \cmd{git add} → \cmd{git commit} してから \cmd{git push} します。
\end{warnbox}

% ---------------------------------------------------------------------
\section{ダウンロードする（pull と fetch）}

別のPCや他の人がリモートを更新したとき、その最新を手元に取り込む（または確認する）のが \cmd{pull} と \cmd{fetch} です。

\begin{center}
\begin{tabularx}{\linewidth}{@{}l X X@{}}
\toprule
コマンド & 動き & いつ使うか \\
\midrule
\cmd{git fetch} & リモートの最新を\textbf{取得して \cmd{origin/main} に記録するだけ}（手元の \cmd{main} や作業ツリーは変えない） & 反映前に「何が来たか」を確認したいとき \\
\cmd{git pull} & \cmd{fetch} + \textbf{統合（merge）}。取得して手元の \cmd{main} にも反映 & ふだんの「最新に追いつく」操作 \\
\bottomrule
\end{tabularx}
\end{center}

\begin{hintbox}
\textbf{\cmd{origin/main} とは？} \cmd{git fetch} で取得したリモートの内容は、手元の \cmd{main} とは別の \cmd{origin/main} という場所に記録されます。これは「\textbf{リモートが今どうなっているかを指す、手元にあるブックマーク}」だと考えてください。だから fetch した後は、自分の作業を変えずに \cmd{git log origin/main} でリモート側の中身を確認できます。
\end{hintbox}

\begin{codebox}[bash]
\begin{code}
# ふだんはこれ（取得して手元の main に反映）
git pull

# 取得だけしておいて、中身を確認してから反映したい場合
git fetch              # リモートの最新を origin/main に取得（手元の main は変わらない）
git log origin/main    # 取得した origin/main の履歴を確認
git merge origin/main  # 問題なければ手元の main に統合
\end{code}
\end{codebox}

\begin{hintbox}
\textbf{作業を始める前に \cmd{git pull} する}のを習慣にすると、手元を常に最新の状態から始められ、後述のコンフリクト（衝突）が起きにくくなります。
\end{hintbox}

\begin{hintbox}
新しめのGit（2.27以降）では、\textbf{初めて \cmd{git pull} したときに一度だけ}「取り込み方を merge にするか rebase にするか」を設定するよう促されることがあります。迷ったら、第1章と同じ標準の \textbf{merge 方式}を選ぶ次の設定を1度だけ実行しておけば、以降は聞かれません。
\begin{codebox}[bash]
\begin{code}
git config --global pull.rebase false
\end{code}
\end{codebox}
これは上流ブランチの設定（\cmd{-u}）とは別の全体設定です。一度決めれば、あとは引数なしの \cmd{git pull} / \cmd{git push} だけで同期できます。
\end{hintbox}

\begin{warnbox}
複数人・複数PCで同じブランチを編集していると、\cmd{pull} 時に\textbf{コンフリクト}（同じ箇所を別々に変更したことによる衝突）が起きることがあります。その解決方法は第3章で詳しく扱います。
\end{warnbox}

% ---------------------------------------------------------------------
\section{日常のフロー}

GitHubとつないだ後の、毎日の基本的な流れです。\textbf{始めにpull、終わりにpush}を挟むだけで、第1章の黄金フロー（編集 → add → commit）がそのまま使えます（図\ref{fig:daily-flow}）。

\begin{numfig}{日常のフロー：pull → 編集 → add → commit → push の繰り返し}{fig:daily-flow}
\begin{tikzpicture}[scale=1.15, transform shape]
  \node[flownode, fill=gitBlue] (P) {git pull\\{\scriptsize 最新を取り込む}};
  \node[flownode, fill=gitRed, right=2.6cm of P] (E) {編集};
  \node[flownode, fill=gitGreen, right=2.6cm of E] (C) {git commit};
  \node[flownode, fill=gitBlue, right=2.6cm of C] (U) {git push\\{\scriptsize クラウドへ反映}};
  \draw[flowarrow] (P) -- (E);
  \draw[flowarrow] (E) -- (C) node[edgelabel, midway, above] {git add};
  \draw[flowarrow] (C) -- (U);
  \draw[flowarrow] (U) to[out=-30, in=-150]
        node[edgelabel, midway, below] {繰り返す} (P);
\end{tikzpicture}
\end{numfig}

\begin{codebox}[bash]
\begin{code}
# 1. まず現在の状態を確認する（今どのファイルがどのエリアにあるか）
git status

# 2. 作業開始前に、GitHub（リモート）の最新を手元に取り込む
git pull

# 3. ファイルを編集したら、変更を記録する（第1章の黄金フロー）
git add .                      # 変更をステージングに追加
git commit -m "変更内容のメモ"   # コミットとして記録

# 4. 記録した変更を GitHub（リモート）へアップロードする
git push
\end{code}
\end{codebox}

この繰り返しで、手元とクラウドが常に同期された状態を保てます。

% ---------------------------------------------------------------------
\section*{この章のまとめ}
\addcontentsline{toc}{section}{この章のまとめ}

\begin{itemize}[leftmargin=1.4em]
  \item \textbf{Git（手元の道具）と GitHub（クラウドのサービス）は別物}。GitHubはローカルリポジトリの「もう一つの倉庫」
  \item リモートとの通信は \textbf{\cmd{clone} / \cmd{push} / \cmd{fetch} / \cmd{pull}} の4つ
  \item \cmd{push} する前に\textbf{認証}（HTTPS + トークン、または SSH鍵）を1度だけ設定する
  \item つなぎ方は2パターン：\textbf{GitHubから始めるなら \cmd{clone}}、\textbf{手元を上げるなら \cmd{remote add} + \cmd{push -u}}
  \item 日常は \textbf{「\cmd{pull} → 編集 → \cmd{add} → \cmd{commit} → \cmd{push}」} の繰り返し
\end{itemize}

次は第3章で、\textbf{ブランチを分けた共同作業・プルリクエスト（PR）・コンフリクトの解決}を扱います。1人で使う分にはここまでで十分ですが、チームで開発するなら第3章が本番です。


\end{document}
